% !TEX root = ../b-rg.tex
\chapter{Prepositions}\label{ch:prepositions}
\stub

% --- Planned sections ---
% \section{Simple prepositions}
%   a, de, ne, con, des, par, por, sur, sou, entr, contr, vars, posc, pre,
%   parmy, lon, lonc, jos, cas, reðr, tras, ant, doutr, deur, jusc ...
%   (~55 in the lexicon).
% \section{Complex prepositions}
%   a cost, a plaç de, a mesur de, dessur de, dessou de, graç a,
%   par mejan de, ne souc de ...
% \section{Contraction with the article}
%   Pointer to (or shared table with) Articles and Demonstratives.
% \section{Expressing semantic roles}
%   Organised by meaning rather than by preposition (Lingua questionnaire
%   2.1.1.4-6): recipient and beneficiary, instrument and means,
%   accompaniment (con, and how it relates to e 'and'), cause, purpose,
%   topic ('about'); location, source, goal and path; location in time
%   (je 'on the day of', d'an 'in the year', durant, posc, pre).
% \section{The core prepositions a, de and ne}
%   Their range of uses; 'than' = a?
% \section{Prepositional phrases}
%   Prepositions with pronouns (which pronoun forms?), with infinitives
%   (por + inf 'in order to') and with ig-clauses; prepositions used
%   without an object, as adverbs; can a preposition be stranded?
